\section{Text Matching and Benchmarks}
\label{sec:text-matching-benchmarks}

\subsection{Text Matching}
\label{sec:text-matching}

Scoring a language model's output against a predetermined reference text is the earliest form of automatic evaluation for generated text. This is a natural language generation problem, not one specific to agents or large language models: any system that produces text needs a way to be scored against a reference. BLEU and ROUGE are the traditional metrics for this problem. BLEU compares a generated translation to one or more reference translations through a weighted average of n-gram precision, matched phrase by phrase against the reference \cite{papineniBleuMethodAutomatic2002}. ROUGE applies the same surface-level comparison to summarization, counting overlapping n-grams, word sequences, or word pairs between a generated summary and human-written reference summaries \cite{linROUGEPackageAutomatic2004}. Both metrics reduce evaluation to how closely the generated text matches a reference text, and this reference-based comparison was shown to correlate poorly with human judgment, particularly on tasks that require creativity and diversity \cite{liuGEvalNLGEvaluation2023}.

\subsection{Benchmarks as Capability Exams}
\label{sec:benchmarks-capability-exams}

A benchmark of a large language model is a set of tasks chosen to represent a capability, together with a standardized procedure for running every model on those tasks and a metric for scoring each one. The result is packaged into a leaderboard so that models can be compared under the same protocol \cite{liangHolisticEvaluationLanguage2023}. Running a model on a task requires an \emph{adaptation} procedure. Adaptation is the method that turns a general-purpose language model into a solver for that one task, for example a specific prompting strategy \cite{liangHolisticEvaluationLanguage2023}. The comparison across models is meaningful only if every model is evaluated on the same scenarios and adapted to each task by the same procedure; otherwise, a difference in score can come from the adaptation method rather than from the model itself \cite{liangHolisticEvaluationLanguage2023}. A benchmark may score an output with text-matching metrics such as BLEU and ROUGE, or apply a different metric when the task requires checking something other than how closely the output resembles a reference string.

MMLU (Massive Multitask Language Understanding) is a benchmark of this kind, on a task where the answer is a single choice. It is a multiple-choice test built from 57 subjects across the humanities, sciences, and other domains, and each item is scored by classification accuracy on the selected answer choice \cite{hendrycksMeasuringMassiveMultitask2021}. The correct answer is one of a predetermined set of options, so the model is scored by whether it selects the option matching the key, not by any comparison of generated text against a reference string.

BFCL (Berkeley Function Calling Leaderboard), another benchmark, illustrates the same point as MMLU for a task where the answer is not text at all: the metric fits the kind of answer the task produces. It evaluates a model's ability to produce function calls and scores them through abstract-syntax-tree sub-string matching, used as a proxy for executing the call: a function call is correct if the function name matches exactly and every parameter value falls within its allowed range \cite{patilBerkeleyFunctionCalling2025}. A function call is a structured object, not a string to be compared word by word, so the benchmark checks its syntactic structure and argument values instead of applying a text-matching metric to it. MMLU and BFCL score two different kinds of output, classification accuracy on a predetermined answer choice and structural correctness on a generated call, and in both cases the metric is chosen to fit what the task is asking.

This benchmark model presupposes headroom: the leaderboard has to separate models by capability, which requires that the tasks not already be solved by every model under test.
